\section*{Chapter 7 --- Component testing through the stub}
\addcontentsline{toc}{section}{Chapter 7 --- Component testing through the stub}

\solution{ex:7:1}{Write the bug, watch the test catch it}
The exercise asks for a code change; the questions are answered here.
\ans{a} The cases checking the echoed frame's content and the counter fail --- the frame that is
echoed has the wrong \code{dlc}, and the case calling \code{run()} twice on one injected frame
reveals that the state has been consumed. The case without an injected frame is still green, since
\code{receive()} returns \code{false} already the first time.
\ans{b} No, neither of them. \code{Stub} and \code{Frame} do exactly what they promise; neither of
them is broken.
\ans{c} A unit test cannot see a bug that lives \emph{between} the units. It takes a test running
several parts together to discover that one of them uses the other incorrectly.

\solution{ex:7:2}{The layer boundary}
\ans{a} No. \code{EchoNode} has to know the interface --- it is its only dependency.
\ans{b} No. The test holds the concrete type; that is where \code{inject()} lives.
\ans{c} Yes. The implementation file is also above the interface, and must not know which driver
it is running against.
\ans{d} No, it does not break the boundary --- \code{hasError()} and \code{clearError()} belong to
the interface. But it is questionable anyway: \code{EchoNode} cannot do anything meaningful about
an error it cannot interpret, and the error flag is one bit for four causes. Error handling
belongs with whoever owns the driver and knows the system's requirements, that is, in the factory
or in the layer above.
\ans{e} Yes, grossly. The cast names a concrete driver \emph{and} reaches a private method in it.
It would also crash as soon as a \code{Stub} was injected.

\solution{ex:7:3}{Ownership}
\ans{a} It does not compile. \code{Interface} is abstract and cannot be instantiated as a value;
the compiler says that the type is abstract. Were it not, the member would \emph{slice} the object
--- copy the base class part and throw the rest away --- which is worse than a compile error.
\ans{b} It compiles. The reference then points at a destroyed object, and \code{run()} is
undefined behaviour. The compiler cannot see it; that is the one cost of the component not owning
its driver.
\ans{c} Because \code{EchoNode} does not own it. The driver's lifetime is longer than the
component's, it can be shared with other components, and the ownership sits with the factory. A
\code{unique_ptr} would have claimed the opposite, and a \code{shared_ptr} would have added
counting that nobody needs.

\solution{ex:7:4}{Read the suite}
\ans{a} The case checking only the echoed frame's \emph{content} --- the right \code{id},
\code{dlc} and data --- would be green, since the echo still happens. That is why the suite also
checks the counter.
\ans{b} Traffic is fed in with \code{stub.inject()} and read back with \code{stub.receive()}.
Neither of them is in \code{Interface}: \code{inject()} \code{Spi} cannot implement, and
\code{receive()} is used here in a different role from the interface's --- the test inspects what
the component \emph{sent}, which the stub happens to store in the same field.
\ans{c} That \code{run()} consumes the frame: after the first call there is nothing left to echo,
so the counter is to stand at 1 and not 2. That is the requirement catching the double call in
exercise~\ref{ex:7:1}.
